% Generated draft: requires mathematical and source review.
\documentclass[11pt]{article}
\usepackage[margin=25mm]{geometry}
\usepackage{fontspec}
\setmainfont{Latin Modern Roman}
\usepackage{amsmath,amssymb,mathtools,unicode-math}
\setmathfont{Latin Modern Math}
\usepackage{enumitem,needspace,xurl,hyperref}
\setlength{\parindent}{0pt}
\setlength{\parskip}{5pt}
\setlength{\emergencystretch}{3em}
\newcommand{\N}{\mathbb{N}}
\newcommand{\Z}{\mathbb{Z}}
\newcommand{\Q}{\mathbb{Q}}
\newcommand{\R}{\mathbb{R}}
\newcommand{\C}{\mathbb{C}}
\newcommand{\F}{\mathbb{F}}
\newcommand{\A}{\mathbb{A}}
\newcommand{\PP}{\mathbb{P}}
\newcommand{\E}{\mathbb{E}}
\newcommand{\eps}{\varepsilon}
\newcommand{\dd}{\,\mathrm{d}}
\begin{document}
\section*{UnsolvedMath ID: 701133}
\textbf{Source problem code:} catalogue entry
\subsection{Definitions and mathematical statement}
Let \(\mathfrak F\) denote the class of all finite groups. For a class \(\mathfrak X\) of groups, define \(\mathbf H_1\mathfrak X\) to consist of groups \(G\) admitting an admissible action on a finite-dimensional contractible complex \(X\) such that every cell stabilizer belongs to \(\mathfrak X\). For precision in this draft, a complex is taken to be a CW complex, and an admissible action is a cellular action for which every element stabilizing a cell setwise fixes that cell pointwise. Thus, writing
\[
G_\sigma=\{g\in G:g\sigma=\sigma\},
\]
the stabilizer requirement is \(G_\sigma\in\mathfrak X\) for every cell \(\sigma\) of \(X\). The supplied record does not explicitly specify the category of complexes or define admissibility, so these conventions require confirmation. The complex need not have finitely many cells; its dimension is finite but may depend on \(G\).

A class \(\mathfrak C\) is \(\mathbf H_1\)-closed if \(\mathbf H_1\mathfrak C\subseteq\mathfrak C\). Define \(\mathbf H\mathfrak X\) to be the smallest \(\mathbf H_1\)-closed class containing \(\mathfrak X\). Groups in \(\mathbf H\mathfrak F\) are the hierarchically decomposable groups considered here. Set
\[
\mathbf H_3\mathfrak F
=\mathbf H_1\bigl(\mathbf H_1(\mathbf H_1\mathfrak F)\bigr).
\]

\textbf{Part (a), historical question, reported solved in the supplied record.} Find a group \(G\) such that
\[
G\in\mathbf H\mathfrak F\setminus\mathbf H_3\mathfrak F.
\]

\textbf{Transfinite indexing qualification.} The source specifies transfinite induction only as being performed ``in the obvious way.'' One conventional completion, used for the conditional formulation below, is
\[
\mathbf H_0\mathfrak F=\mathfrak F,\qquad
\mathbf H_\alpha\mathfrak F
=\mathbf H_1\left(\bigcup_{\beta<\alpha}\mathbf H_\beta\mathfrak F\right)
\quad(\alpha>0).
\]
Here each cell stabilizer belongs to some earlier stage, and the earlier ordinal may depend on the cell; a single earlier stage containing all stabilizers is not required by this formula. In particular, this recursion applies \(\mathbf H_1\) at limit ordinals as well. The supplied text does not explicitly distinguish this convention from a recursion that takes only the union of earlier stages at limit ordinals. The intended limit-stage definition must therefore be confirmed.

\textbf{Part (b), conditional on the intended transfinite convention.} Determine whether there exists an ordinal \(\alpha\) such that
\[
\mathbf H_\alpha\mathfrak F=\mathbf H\mathfrak F,
\]
meaning that, for every group \(G\),
\[
G\in\mathbf H_\alpha\mathfrak F
\quad\Longleftrightarrow\quad
G\in\mathbf H\mathfrak F.
\]
The quantifier over \(\alpha\) is not restricted to countable ordinals, and no countability restriction on \(G\) is supplied.

\textbf{Supplied status and qualifications.} The registry labels the overall problem partially solved. It reports part (a) solved and attributes to Januszkiewicz, Kropholler, and Leary the assertion that, for every countable ordinal \(\alpha\), there is a group in
\(\mathbf H_{\alpha+1}\mathfrak F\setminus\mathbf H_\alpha\mathfrak F\), with the hierarchy understood in the source's intended sense. The editor consequently reports that an ordinal satisfying part (b) cannot be countable. The background also reports torsion-free examples with similar properties, attributed to Fournier-Facio and Sun, but does not specify their exact quantifiers or hierarchy properties. These are supplied literature assessments, not independently verified conclusions of this draft.
\subsection{Short English statement}
The problem asks how many stages of actions on finite-dimensional contractible complexes are needed to obtain all hierarchically decomposable groups. Part (a), reported solved, asks for a group beyond the third stage. Part (b) asks whether any ordinal stage contains the entire class; the supplied source reports that no countable stage does. The precise limit-stage recursion and action conventions require confirmation.
\subsection{Sources}
\begin{itemize}
\item Supplied local UnsolvedMath registry record, permanent ID \(701133\), exact source problem code \texttt{catalogue entry}, titled \emph{Kourovka Notebook Problem 15.45}. Proposed by P. Kropholler in 2002. The supplied registry status is \texttt{partially\_solved}.
\item \emph{The Kourovka Notebook (Unsolved Problems in Group Theory)}, original problem 15.45, 15th issue (2002); supplied updated source: 21st issue edition (2026), arXiv:\texttt{1401.0300v48}. The collection description names E. I. Khukhro and V. D. Mazurov as editors. The supplied status-review note identifies the edition as October 2026.
\item T. Januszkiewicz, P. H. Kropholler, and I. J. Leary, \emph{Bull. London Math. Soc.} \textbf{42}, no. 5 (2010), 896--904. The supplied background attributes the solution of part (a) and strictness at every countable successor stage to this reference; the editor's comment attributes the exclusion of a countable ordinal in part (b) to it. No article title is supplied.
\item F. Fournier-Facio and B. Sun, preprint (2025), \url{https://arxiv.org/abs/2503.01987v3}. The supplied background describes torsion-free examples with similar properties, without specifying a precise theorem or article title.
\item The supplied status review has exact checking date \texttt{2026-10-16} and evidence designation \texttt{maintained\_source}. Its partial-solution assessment is supplied dated curation, not independently verified here. No cited source was independently consulted for this draft.
\end{itemize}
\end{document}
